\section{Ipergrafi Temporali}\label{sec:ipergrafi_temporali}

Come abbiamo descritto in precedenza gli ipergrafi statici possono risultare la migliore struttura matematica da utilizzare per la rappresentazione di strutture complesse senza causare una perdita di informazioni, però tale struttura non può essere applicare direttamente a contesti applicativi senza prima 



% prima versione
\section{Ipergrafi Temporali}\label{sec:ipergrafi_temporali}

Nonostante gli ipergrafi consentano di rappresentare interazioni che coinvolgono più entità contemporaneamente, la formulazione classica presentata finora descrive esclusivamente la struttura statica del sistema, ignorando la dimensione temporale delle interazioni.
Il vero potenziale degli ipergrafi emerge quando vengono considerati interazioni reali e complesse, come quelle che si verificano tra individui all'interno di una rete sociale.
All'interno di contesti simili, a causa della natura dinamica delle interazioni stesse, l'utilizzo di una rappresentazione statica dell'ipergrafo non è sufficiente per catturare la complessità delle relazioni tra gli individui e può comportare una perdita di informazioni relative alla loro evoluzione temporale.
Per questo motivo, è necessario estendere il concetto di ipergrafo integrando la dimensione temporale.
La struttura dati risultante prende il nome di \textit{ipergrafo temporale} e consente di rappresentare simultaneamente la natura collettiva delle interazioni e la loro evoluzione nel tempo~\cite{cencetti2021}.
Formalmente, un ipergrafo temporale è una coppia $H_t = (V, E_T)$, dove $V$ è l'insieme di nodi e $E_T$ è l'insieme di coppie $(E, t)$, definito come:
\[
    E_T \subseteq \{(E,t) \mid E \subseteq V, E \neq \emptyset, t \in T\}.
\]
Ogni elemento $(E,t)$ rappresenta un'interazione che coinvolge simultaneamente tutti i nodi appartenenti all'iperarco $E$ all'istante temporale $t$.
Genericamente parlando, a seconda del contesto applicativo, l'informazione temporale può essere associata a un singolo istante oppure ad un intervallo di tempo.
L'introduzione della componente temporale consente di distinguere tra sistemi che, pur presentando la stessa struttura aggregata, differiscono per l'ordine cronologico delle interazioni.
Tale aspetto risulta particolarmente rilevante nello studio di processi dinamici, nei quali la sequenza temporale degli eventi può influenzare significativamente l'evoluzione del sistema~\cite{holme2012}.

\subsection{Finestra temporale}

In molti contesti applicativi, considerare tutte le interazioni avvenute dall'inizio del periodo di osservazione può risultare poco significativo, poiché le relazioni più recenti sono spesso più rilevanti rispetto a quelle avvenute in passato.
Per questo motivo, nell'analisi degli ipergrafi temporali viene introdotto il concetto di \textit{finestra temporale}.
Data una durata $(\Delta T)$, ogni interazione osservata all'istante $t$ viene considerata valida solamente nell'intervallo temporale $[t, t + \Delta T]$.
Trascorso tale intervallo, l'interazione non viene più considerata nell'analisi del sistema.
In altre parole, dato un generico istante $t_c$, l'ipergrafo è costituito esclusivamente dagli iperarchi associati a eventi verificatisi nell'intervallo $[t_c - \Delta T, t_c]$.
Questo approccio consente di focalizzare l'attenzione sulle interazioni più recenti, trascurando progressivamente quelle meno rilevanti.
Ad esempio, in un social network può essere utile analizzare solamente le conversazioni avvenute negli ultimi sette giorni, mentre in un contesto aziendale si potrebbe considerare un orizzonte temporale pari a un mese.
La finestra temporale non rimane necessariamente fissa durante l'analisi.
A partire da una posizione iniziale, essa viene fatta scorrere lungo l'asse temporale mediante un passo di avanzamento $\delta$, che può essere differente dalla sua ampiezza $\Delta T$. Ad ogni avanzamento viene costruita una nuova istanza dell'ipergrafo contenente esclusivamente le interazioni comprese nella nuova finestra di osservazione.
La scelta dei parametri $\Delta T$ e $\delta$ influenza significativamente le proprietà della rappresentazione ottenuta e vengono fornite dall'utente in base al contesto applicativo.
Finestre temporali più ampie tendono a produrre ipergrafi più densi, poiché includono un numero maggiore di interazioni, mentre finestre più ristrette permettono di descrivere con maggiore precisione l'evoluzione temporale del sistema.
Analogamente, un passo di avanzamento ridotto consente di osservare variazioni più graduali della struttura dell'ipergrafo, mentre valori più elevati producono una descrizione temporale più grossolana ma computazionalmente meno onerosa.

